\chapter{Thực nghiệm và kết quả}
\section{Thiết lập thực nghiệm}
Thực nghiệm sentiment sử dụng dữ liệu 5.162 review sau dedup. Dữ liệu được chia train/valid/test với phân phối lớp như Chương 4. Các mô hình so sánh gồm majority baseline, \TFIDF{} + Logistic Regression và \TFIDF{} + Linear SVM. Tiêu chí chọn model là validation \MacroFOne{}, vì \MacroFOne{} đánh giá đều các lớp và phù hợp khi lớp neutral ít hơn.

Thực nghiệm retrieval/RAG dùng 15 query đánh giá trong Module 4. Kết quả retrieval được đo bằng score trung bình tự động và manual Precision@k. Score embedding không đủ để kết luận relevance, vì vậy Precision@k thủ công là số liệu quan trọng hơn. Module 6 dùng benchmark nhỏ với keyword hit-rate và aspect hit-rate để so sánh BM25, dense, hybrid và hybrid + reranker.

\section{Kết quả sentiment classification}
\begin{table}[H]\centering\caption{So sánh mô hình sentiment trên tập test}\label{tab:sentiment-model-comparison}\begin{tabular}{lrrr}\toprule Mô hình & Accuracy & Macro-F1 & Weighted-F1 \\\midrule Majority baseline & 0,4297 & 0,2004 & 0,2583 \\ \TFIDF{} + Logistic Regression & 0,9174 & 0,8940 & 0,9176 \\ \TFIDF{} + Linear SVM & \textbf{0,9329} & \textbf{0,9098} & \textbf{0,9315} \\\bottomrule\end{tabular}\end{table}

Bảng \ref{tab:sentiment-model-comparison} cho thấy hai mô hình \TFIDF{} đều vượt xa majority baseline. Linear SVM là mô hình tốt nhất trong các baseline đã thử. So với majority baseline, Linear SVM tăng test accuracy từ 0,4297 lên 0,9329 và \MacroFOne{} từ 0,2004 lên 0,9098. Điều này chứng minh mô hình thật sự học được tín hiệu từ văn bản, không chỉ dựa vào lớp nhiều nhất.

\safeimage{figures/module2/model_comparison_accuracy.png}{So sánh accuracy giữa các mô hình sentiment.}{fig:model-accuracy}{0.82\textwidth}
\safeimage{figures/module2/model_comparison_macro_f1.png}{So sánh \MacroFOne{} giữa các mô hình sentiment.}{fig:model-macro-f1}{0.82\textwidth}

\section{Per-class F1 của Linear SVM}
\begin{table}[H]\centering\caption{Per-class metric của Linear SVM trên tập test}\label{tab:linear-svm-per-class}\begin{tabular}{lrrrr}\toprule Lớp & Precision & Recall & F1-score & Support \\\midrule Negative & 0,9102 & 0,9775 & 0,9426 & 311 \\ Neutral & 0,8929 & 0,7634 & 0,8230 & 131 \\ Positive & 0,9696 & 0,9580 & 0,9637 & 333 \\ Macro avg & 0,9242 & 0,8996 & 0,9098 & 775 \\ Weighted avg & 0,9328 & 0,9329 & 0,9315 & 775 \\\bottomrule\end{tabular}\end{table}

Lớp positive và negative có F1 cao hơn neutral. Đây là hiện tượng hợp lý vì hai cực cảm xúc thường có từ khóa rõ như “tốt”, “hài lòng”, “tệ”, “lỗi”, “không trả lời”. Lớp neutral ít dữ liệu hơn và thường chứa nhận xét lưng chừng hoặc mixed sentiment, nên dễ bị nhầm.


\begin{table}[H]
\centering
\caption[Confusion matrix dạng số của Linear SVM]{Confusion matrix dạng số của Linear SVM trên tập test}
\label{tab:linear-svm-confusion-counts}
\begin{tabular}{lrrr}
\toprule
True / Predicted & Negative & Neutral & Positive \\
\midrule
Negative & \textbf{304} & 5 & 2 \\
Neutral & 23 & \textbf{100} & 8 \\
Positive & 7 & 7 & \textbf{319} \\
\bottomrule
\end{tabular}
{\footnotesize Nguồn: \filepath{results/module2/linear_svm_predictions.csv}.}
\end{table}
Từ Bảng \ref{tab:linear-svm-confusion-counts}, negative đúng 304/311 và positive đúng 319/333. Neutral đúng 100/131, trong đó 23 review neutral bị nhầm thành negative. Đây là bằng chứng định lượng cho nhận xét rằng neutral là lớp khó nhất.

\safeimage{figures/module2/linear_svm_confusion_matrix.png}{Confusion matrix của Linear SVM.}{fig:svm-confusion}{0.72\textwidth}
\safeimage{figures/module2/logistic_regression_confusion_matrix.png}{Confusion matrix của Logistic Regression.}{fig:logreg-confusion}{0.72\textwidth}

\section{Kết quả retrieval/RAG Module 4}
Manual Precision@k của retrieval baseline được tổng hợp từ kết quả đánh giá thủ công của Module 5. Kết quả gồm 15 query đánh giá.
\begin{table}[H]\centering\caption{Manual Precision@k của retrieval baseline}\label{tab:manual-precision}\begin{tabular}{lrrr}\toprule Metric & Giá trị & Số query & Số dòng đã chấm \\\midrule Precision@3 & 0,4222 & 15 & 45 \\ Precision@5 & 0,5067 & 15 & 75 \\ Precision@10 & 0,4867 & 15 & 150 \\\bottomrule\end{tabular}\end{table}

Precision@5 khoảng 0,5067 nghĩa là trung bình trong 5 review đầu được trả về, khoảng một nửa được chấm là liên quan. Đây là mức baseline, chưa phải retrieval hoàn hảo. Precision@5 cao hơn Precision@3 trong bộ đánh giá hiện có có thể do một số query có evidence liên quan xuất hiện ở rank 4--5; với tập query nhỏ, biến động này là bình thường.

\safeimage{figures/module4/retrieval_eval_summary.png}{Tổng hợp đánh giá retrieval Module 4.}{fig:retrieval-summary}{0.82\textwidth}
\safeimage{figures/module5/manual_precision_at_k.png}{Manual Precision@k của retrieval baseline.}{fig:manual-pk}{0.78\textwidth}

\section{Score retrieval tự động}
Theo summary retrieval của Module 4, average retrieval score trung bình xấp xỉ 0,8420 với k=3, 0,8404 với k=5 và 0,8381 với k=10. Score giảm nhẹ khi k tăng là hợp lý vì các kết quả sau thường ít gần query hơn. Tuy nhiên, score embedding không phải nhãn relevance. Một review có score cao vẫn có thể không trả lời đúng intent nếu query mơ hồ hoặc review cùng chủ đề nhưng khác polarity.
\begin{table}[H]\centering\caption{Average retrieval score theo k}\label{tab:retrieval-score}\begin{tabular}{lr}\toprule k & Average score xấp xỉ \\\midrule 3 & 0,8420 \\ 5 & 0,8404 \\ 10 & 0,8381 \\\bottomrule\end{tabular}\end{table}

\section{Kết quả Module 6}
Module 6 được đánh giá bằng benchmark nhỏ kiểm tra keyword hit-rate và aspect hit-rate. Benchmark này không thay thế đánh giá relevance thủ công lớn.
\begin{table}[H]\centering\caption{So sánh Module 6 trên benchmark nhỏ}\label{tab:module6-comparison}\begin{tabular}{lrrrr}\toprule Method & keyword\_hit@5 & keyword\_hit@10 & aspect\_hit@5 & aspect\_hit@10 \\\midrule BM25 & 0,500 & 0,667 & 1,000 & 1,000 \\ Dense & 0,500 & 0,833 & 1,000 & 1,000 \\ Hybrid & 1,000 & 1,000 & 1,000 & 1,000 \\ Hybrid + reranker & 1,000 & 1,000 & 1,000 & 1,000 \\\bottomrule\end{tabular}\end{table}

Kết quả cho thấy hybrid retrieval cải thiện khả năng bắt keyword trong benchmark nhỏ. Điều này phù hợp với giả thuyết: BM25 mạnh ở exact keyword, dense mạnh ở semantic similarity, còn hybrid kết hợp hai ưu điểm. Tuy nhiên, do benchmark nhỏ và handcrafted, kết quả này chỉ nên xem là tín hiệu kỹ thuật ban đầu.
\safeimage{figures/module6/module6_retrieval_comparison_reference.png}{So sánh BM25, Dense, Hybrid và Hybrid + reranker trong Module 6.}{fig:module6-comparison}{0.82\textwidth}

\section{Nhận xét tổng hợp}
Sentiment classification là phần đạt kết quả mạnh nhất, với Linear SVM \MacroFOne{} khoảng 0,9098. Retrieval/RAG đã hoạt động end-to-end và có citation, nhưng Precision@5 còn ở mức baseline khoảng 0,5067. Module 6 chỉ ra hướng cải thiện retrieval theo aspect, nhưng cần đánh giá lớn hơn trước khi khẳng định hiệu quả tổng quát.

\section{Diễn giải sâu kết quả sentiment}
Linear SVM vượt Logistic Regression khoảng 0,0158 \MacroFOne{} trên tập test. Chênh lệch này không quá lớn nhưng nhất quán với validation \MacroFOne{}. Điều này cho thấy Linear SVM phù hợp hơn với dữ liệu \TFIDF{} sparse trong đề tài. Majority baseline có \MacroFOne{} rất thấp vì chỉ dự đoán lớp positive, làm F1 của negative và neutral bằng 0.

Lớp positive có F1 cao nhất, khoảng 0,9637. Lớp negative cũng cao, khoảng 0,9426. Điều này cho thấy từ khóa tích cực và tiêu cực trong review khá rõ. Lớp neutral có F1 khoảng 0,8230, thấp hơn rõ rệt. Đây là dấu hiệu quan trọng: nếu muốn cải thiện model, nên tập trung vào neutral/mixed, không chỉ tối ưu accuracy tổng thể.

\section{Diễn giải sâu kết quả retrieval}
Manual Precision@5 bằng 0,5067 không phải con số thấp tuyệt đối nếu xét đây là baseline dense retrieval trên review nhiễu, nhưng cũng chưa đủ để gọi là hệ thống retrieval mạnh. Hệ thống có thể trả kết quả phù hợp hơn với query cụ thể như giao hàng chậm, sai hàng, thiếu hàng, đóng gói lỗi. Với query quá rộng, top-k dễ lẫn nhiều aspect.

Precision@10 thấp hơn Precision@5 một chút vì khi lấy nhiều kết quả hơn, các rank sau thường ít liên quan hơn. Đây là hiện tượng phổ biến trong retrieval. Nếu dùng RAG answer, top-k quá lớn có thể đưa thêm noise vào context. Vì vậy giao diện mặc định top-k = 5 là hợp lý: đủ evidence để tổng hợp nhưng không quá nhiều để nhiễu.

\section{Giới hạn của benchmark Module 6}
Module 6 benchmark dùng các query và keyword mục tiêu được thiết kế thủ công. Chỉ số \code{keyword_hit@5} trả lời câu hỏi: trong top-5 có review chứa keyword mục tiêu không. Chỉ số \code{aspect_hit@5} trả lời câu hỏi: trong top-5 có review thuộc aspect mục tiêu không. Đây là proxy metric, không phải relevance label đầy đủ.

Vì vậy, việc Hybrid đạt 1,000 trên benchmark nhỏ cần diễn giải thận trọng. Nó cho thấy hybrid retrieval bắt keyword/aspect tốt hơn trong các case đã thiết kế, nhưng chưa chứng minh mọi query thực tế đều tốt hơn. Muốn kết luận chắc chắn, cần mở rộng benchmark, chấm relevance thủ công và đo Precision@k/NDCG/MRR trên tập query lớn hơn.

\section{Threats to validity}
Có ba nhóm đe dọa tính hợp lệ. Thứ nhất là dữ liệu: nhãn sentiment lấy từ rating nên có thể nhiễu. Thứ hai là đánh giá: retrieval manual chỉ có 15 query, chưa đại diện hết mọi nhu cầu. Thứ ba là triển khai: kết quả giao diện thử nghiệm phụ thuộc artifact hiện có, nếu thay corpus hoặc embedding model thì cần build lại index và đánh giá lại. Báo cáo đã ghi rõ các điểm này để tránh phóng đại.

\section{So sánh kết quả với kỳ vọng ban đầu}
Trước khi chạy thực nghiệm, có thể kỳ vọng Logistic Regression và Linear SVM đều mạnh vì dữ liệu review ngắn và có nhiều từ khóa sentiment. Kết quả thực tế đúng với kỳ vọng này: cả hai mô hình đều vượt majority baseline rất xa. Linear SVM tốt hơn Logistic Regression, nhưng Logistic Regression vẫn đạt \MacroFOne{} 0,8940, chứng tỏ \TFIDF{} là biểu diễn phù hợp cho baseline.

Kết quả retrieval lại khiêm tốn hơn. Điều này cũng hợp lý vì retrieval khó hơn classification trong đề tài này. Classification chỉ chọn một trong ba lớp đã định nghĩa. Retrieval phải trả về đúng review liên quan cho nhiều query mở. Một query có thể liên quan đến nhiều aspect, nhiều cách diễn đạt và nhiều mức sentiment. Vì vậy Precision@5 khoảng 0,5067 là kết quả baseline cần được cải thiện, không phải thất bại của hệ thống.

\section{Vai trò của trực quan hóa trong phân tích kết quả}
Các biểu đồ accuracy và \MacroFOne{} giúp so sánh nhanh hiệu quả giữa các mô hình sentiment. Trong đó, \MacroFOne{} có ý nghĩa quan trọng hơn accuracy vì phản ánh hiệu quả trên cả lớp neutral, vốn là lớp ít mẫu và khó phân biệt hơn. Confusion matrix bổ sung góc nhìn chi tiết hơn bằng cách cho thấy mẫu nào bị nhầm giữa negative, neutral và positive.

Đối với retrieval, biểu đồ Precision@k giúp diễn giải mức độ liên quan của các review được truy xuất ở các ngưỡng top-k khác nhau. Trực quan hóa không thay thế bảng số liệu, nhưng hỗ trợ nhận diện xu hướng: sentiment classification đạt kết quả ổn định hơn retrieval, trong khi retrieval vẫn còn dư địa cải thiện ở các query rộng hoặc nhiều aspect.

\section{Diễn giải chỉ số Precision@k}
Với mỗi câu hỏi, hệ thống trả về danh sách review ở top-k. Precision@k là tỷ lệ review liên quan trong k review đầu, trung bình qua các query. Giá trị Precision@5 = 0,5067 nghĩa là trong phạm vi bộ đánh giá 15 query, trung bình khoảng một nửa số review ở top-5 được gán nhãn liên quan.

Đây là baseline manual evaluation, không nên diễn giải thành “độ chính xác RAG là 50\%”. RAG gồm nhiều phần: retrieval, answer formatting và evidence display. Precision@5 chỉ đo relevance của phần retrieval ở top-5. Answer có thể vẫn hữu ích nếu evidence liên quan nằm trong top-5 và hệ thống tổng hợp thận trọng.

\section{Diễn giải kết quả benchmark Module 6}
Module 6 được xem như một hướng cải tiến retrieval. Dense retrieval nắm bắt tương đồng ngữ nghĩa nhưng có thể bỏ sót keyword complaint. BM25 bắt keyword tốt hơn nhưng không hiểu nghĩa rộng. Hybrid dùng cả hai và gộp bằng RRF. Benchmark nhỏ cho thấy hybrid bắt keyword/aspect tốt hơn trong các query thử nghiệm. Tuy nhiên, benchmark còn nhỏ và cần đánh giá lớn hơn trước khi kết luận tổng quát.

